\documentclass[11pt]{article}

\usepackage[utf8]{inputenc}
\usepackage[T1]{fontenc}
\usepackage[spanish]{babel}
\usepackage{amsmath,amssymb,mathtools,bm,graphicx,xcolor}
\usepackage{float}
\usepackage{svg}
\svgsetup{inkscapelatex=false}
\usepackage[a4paper,margin=2.5cm]{geometry}
\usepackage[
    colorlinks=true,
    linkcolor=red,
    citecolor=green,
    urlcolor=red
]{hyperref}

\spanishdecimal{.}

\setlength{\parindent}{0pt}
\setlength{\parskip}{0.45em}
\setlength{\abovedisplayskip}{6pt}
\setlength{\belowdisplayskip}{6pt}
\setlength{\abovedisplayshortskip}{4pt}
\setlength{\belowdisplayshortskip}{4pt}

% Formato de encabezado y problemas
\newcommand{\encabezado}[3]{%
{\Large\bfseries #1\par}
\vspace{2pt}
{\normalsize #2\hfill #3\par}
\vspace{6pt}\hrule\vspace{8pt}}

\newcommand{\problema}[2]{%
\vspace{0.55em}
\noindent\textbf{#1.}\enspace{\small #2}\par
\vspace{0.3em}}

\newcommand{\inciso}[1]{%
\vspace{0.35em}
\noindent\textbf{(#1)}\enspace}

% Comandos auxiliares
\newcommand{\dd}{\,\mathrm{d}}
\newcommand{\ii}{\mathrm{i}}
\newcommand{\e}{\mathrm{e}}
\newcommand{\sen}{\operatorname{sen}}
\newcommand{\grad}{\nabla}
\newcommand{\laplaciano}{\nabla^2}
\newcommand{\R}{\mathbb{R}}
\newcommand{\Z}{\mathbb{Z}}
\newcommand{\N}{\mathbb{N}}
\newcommand{\Nzero}{\mathbb{N}_0}
\newcommand{\JR}{J.R.}
\newcommand{\GR}{G.R.}

\begin{document}

\encabezado{Helmholtz en coordenadas esféricas}{Física Matemática II}{J.R., G.R.}

\problema{1}{Planteamiento del problema}

Consideremos la ecuación de Helmholtz modificada
\begin{equation}
\nabla^2\Psi+\alpha\Psi=0,
    \label{eq:helmholtz}
\end{equation}
donde $\Psi=\Psi(r,\theta,\phi)$ es una función escalar y $\alpha\in\R$ es una constante. En coordenadas esféricas, el Laplaciano se escribe como
\begin{equation*}
\nabla^2\Psi
    =
    \frac{1}{r^2}\frac{\partial}{\partial r}
    \left(r^2\frac{\partial\Psi}{\partial r}\right)
    +
    \frac{1}{r^2\sin\theta}\frac{\partial}{\partial\theta}
    \left(\sin\theta\frac{\partial\Psi}{\partial\theta}\right)
    +
    \frac{1}{r^2\sin^2\theta}\frac{\partial^2\Psi}{\partial\phi^2}.
\end{equation*}
Por lo tanto, la ecuación \eqref{eq:helmholtz} adopta la forma explícita
\begin{equation}
\frac{1}{r^2}\frac{\partial}{\partial r}
    \left(r^2\frac{\partial\Psi}{\partial r}\right)
    +
    \frac{1}{r^2\sin\theta}\frac{\partial}{\partial\theta}
    \left(\sin\theta\frac{\partial\Psi}{\partial\theta}\right)
    +
    \frac{1}{r^2\sin^2\theta}\frac{\partial^2\Psi}{\partial\phi^2}
    +\alpha\Psi=0.
    \label{eq:helmholtz_esferica}
\end{equation}
El objetivo es aplicar el método de separación de variables hasta obtener las EDOs para las partes radial, polar y azimutal.

\problema{2}{Ansatz separable}

Buscaremos soluciones separables de la forma
\begin{equation}
\Psi_{\mathrm{sep}}(r,\theta,\phi)=R(r)\Theta(\theta)\Phi(\phi).
    \label{eq:ansatz_helmholtz}
\end{equation}
Este ansatz separa la dependencia radial, polar y azimutal. Si una solución de la forma \eqref{eq:ansatz_helmholtz} satisface la EDP, entonces por linealidad podremos construir soluciones más generales como superposiciones de soluciones separables.

\problema{3}{Sustitución explícita en la EDP}

Calculamos primero las derivadas parciales que aparecen en \eqref{eq:helmholtz_esferica}. Como $\Theta$ y $\Phi$ no dependen de $r$,
\begin{equation*}
\frac{\partial\Psi_{\mathrm{sep}}}{\partial r}
    =
    \Theta(\theta)\Phi(\phi)\frac{\dd R}{\dd r}(r),
\end{equation*}
y entonces
\begin{equation}
\frac{\partial}{\partial r}
    \left(r^2\frac{\partial\Psi_{\mathrm{sep}}}{\partial r}\right)
    =
    \Theta(\theta)\Phi(\phi)
    \frac{\dd}{\dd r}\left(r^2\frac{\dd R}{\dd r}\right).
    \label{eq:derivada_r_2}
\end{equation}

Análogamente, como $R$ y $\Phi$ no dependen de $\theta$,
\begin{equation*}
\frac{\partial\Psi_{\mathrm{sep}}}{\partial\theta}
    =
    R(r)\Phi(\phi)\frac{\dd\Theta}{\dd\theta}(\theta),
\end{equation*}
y por lo tanto
\begin{equation}
\frac{\partial}{\partial\theta}
    \left(\sin\theta\frac{\partial\Psi_{\mathrm{sep}}}{\partial\theta}\right)
    =
    R(r)\Phi(\phi)
    \frac{\dd}{\dd\theta}
    \left(\sin\theta\frac{\dd\Theta}{\dd\theta}\right).
    \label{eq:derivada_theta_2}
\end{equation}
Finalmente, como $R$ y $\Theta$ no dependen de $\phi$,
\begin{equation}
\frac{\partial^2\Psi_{\mathrm{sep}}}{\partial\phi^2}
    =
    R(r)\Theta(\theta)\frac{\dd^2\Phi}{\dd\phi^2}(\phi).
    \label{eq:derivada_phi_2}
\end{equation}

Sustituyendo \eqref{eq:derivada_r_2}, \eqref{eq:derivada_theta_2} y \eqref{eq:derivada_phi_2} en \eqref{eq:helmholtz_esferica}, encontramos
\begin{align*}
&\frac{1}{r^2}\Theta\Phi
    \frac{\dd}{\dd r}\left(r^2\frac{\dd R}{\dd r}\right)
    +
    \frac{1}{r^2\sin\theta}R\Phi
    \frac{\dd}{\dd\theta}\left(\sin\theta\frac{\dd\Theta}{\dd\theta}\right)
    \notag\\
    &\hspace{5em}
    +
    \frac{1}{r^2\sin^2\theta}R\Theta
    \frac{\dd^2\Phi}{\dd\phi^2}
    +
    \alpha R\Theta\Phi=0.
\end{align*}
Dividiendo por $\Psi_{\mathrm{sep}}=R\Theta\Phi$, obtenemos
\begin{equation*}
\frac{1}{r^2R}\frac{\dd}{\dd r}\left(r^2\frac{\dd R}{\dd r}\right)
    +
    \frac{1}{r^2\Theta\sin\theta}
    \frac{\dd}{\dd\theta}\left(\sin\theta\frac{\dd\Theta}{\dd\theta}\right)
    +
    \frac{1}{r^2\sin^2\theta\,\Phi}\frac{\dd^2\Phi}{\dd\phi^2}
    +\alpha=0.
\end{equation*}
Multiplicamos ahora por $r^2\sin^2\theta$. De este modo,
\begin{equation}
\sin^2\theta
    \left[
    \frac{1}{R}\frac{\dd}{\dd r}\left(r^2\frac{\dd R}{\dd r}\right)
    +
    \frac{1}{\Theta\sin\theta}
    \frac{\dd}{\dd\theta}\left(\sin\theta\frac{\dd\Theta}{\dd\theta}\right)
    +
    \alpha r^2
    \right]
    +
    \frac{1}{\Phi}\frac{\dd^2\Phi}{\dd\phi^2}=0.
    \label{eq:helmholtz_forma_separable}
\end{equation}
Esta es la forma conveniente para realizar la primera separación.

\problema{4}{Primera separación: variable azimutal}

En \eqref{eq:helmholtz_forma_separable}, el término
\begin{equation*}
\frac{1}{\Phi}\frac{\dd^2\Phi}{\dd\phi^2}
\end{equation*}
depende sólo de $\phi$, mientras que el resto depende sólo de $r$ y $\theta$. Por lo tanto, ambos términos deben ser constantes. Siguiendo la convención usual, escribimos
\begin{equation}
\frac{1}{\Phi}\frac{\dd^2\Phi}{\dd\phi^2}=\text{const.}=-m^2.
    \label{eq:separacion_phi}
\end{equation}
Así, la EDO azimutal es
\begin{equation*}
\frac{\dd^2\Phi}{\dd\phi^2}+m^2\Phi=0.
\end{equation*}
Una base compleja de soluciones está dada por
\begin{equation}
\Phi(\phi)=\e^{\ii m\phi}.
    \label{eq:solucion_phi_exp}
\end{equation}
Como $\phi$ es una coordenada angular, la solución física debe ser univaluada, es decir,
\begin{equation}
\Phi(\phi+2\pi)=\Phi(\phi).
    \label{eq:periodicidad_phi}
\end{equation}
Sustituyendo \eqref{eq:solucion_phi_exp} en \eqref{eq:periodicidad_phi}, se obtiene
\begin{equation*}
\e^{\ii m(\phi+2\pi)}=\e^{\ii m\phi}
    \quad\Longrightarrow\quad
    \e^{2\pi\ii m}=1.
\end{equation*}
Por consiguiente,
\begin{equation*}
m=0,\pm1,\pm2,\ldots,
    \qquad m\in\Z.
\end{equation*}

\problema{5}{Segunda separación: variables radial y polar}

Sustituyendo \eqref{eq:separacion_phi} en \eqref{eq:helmholtz_forma_separable}, se obtiene
\begin{equation*}
\sin^2\theta
    \left[
    \frac{1}{R}\frac{\dd}{\dd r}\left(r^2\frac{\dd R}{\dd r}\right)
    +
    \frac{1}{\Theta\sin\theta}
    \frac{\dd}{\dd\theta}\left(\sin\theta\frac{\dd\Theta}{\dd\theta}\right)
    +
    \alpha r^2
    \right]
    -m^2=0.
\end{equation*}
Dividiendo por $\sin^2\theta$,
\begin{equation*}
\frac{1}{R}\frac{\dd}{\dd r}\left(r^2\frac{\dd R}{\dd r}\right)
    +
    \alpha r^2
    +
    \frac{1}{\Theta\sin\theta}
    \frac{\dd}{\dd\theta}\left(\sin\theta\frac{\dd\Theta}{\dd\theta}\right)
    -
    \frac{m^2}{\sin^2\theta}=0.
\end{equation*}
Los dos primeros términos dependen sólo de $r$, mientras que los dos últimos dependen sólo de $\theta$. Por lo tanto, introducimos una segunda constante de separación $Q$:
\begin{equation}
\frac{1}{R}\frac{\dd}{\dd r}\left(r^2\frac{\dd R}{\dd r}\right)
    +
    \alpha r^2=Q,
    \label{eq:separacion_radial_Q}
\end{equation}
\begin{equation}
\frac{1}{\Theta\sin\theta}
    \frac{\dd}{\dd\theta}\left(\sin\theta\frac{\dd\Theta}{\dd\theta}\right)
    -
    \frac{m^2}{\sin^2\theta}=-Q.
    \label{eq:separacion_polar_Q}
\end{equation}
Multiplicando \eqref{eq:separacion_radial_Q} por $R$ y \eqref{eq:separacion_polar_Q} por $\Theta$, obtenemos las EDOs separadas:
\begin{align}
\frac{\dd}{\dd r}\left(r^2\frac{\dd R}{\dd r}\right)
    +
    (\alpha r^2-Q)R&=0,
    \label{eq:edo_radial}\\
    \frac{1}{\sin\theta}\frac{\dd}{\dd\theta}
    \left(\sin\theta\frac{\dd\Theta}{\dd\theta}\right)
    +
    \left(Q-\frac{m^2}{\sin^2\theta}\right)\Theta&=0,
    \label{eq:edo_polar}\\
    \frac{\dd^2\Phi}{\dd\phi^2}+m^2\Phi&=0.
    \notag
\end{align}
Con esto queda realizada la separación de variables de la ecuación de Helmholtz en coordenadas esféricas.

\problema{6}{Relación con la ecuación asociada de Legendre}

La EDO polar \eqref{eq:edo_polar} puede escribirse como
\begin{equation}
\frac{1}{\sin\theta}\frac{\dd}{\dd\theta}
    \left(\sin\theta\frac{\dd\Theta}{\dd\theta}\right)
    +
    \left(Q-\frac{m^2}{\sin^2\theta}\right)\Theta=0.
    \label{eq:polar_antes_cambio}
\end{equation}
Realizamos el cambio de variable
\begin{equation*}
x:=\cos\theta,
    \qquad
    \Theta(\theta)=y(x)=y(\cos\theta).
\end{equation*}
Como
\begin{equation*}
\frac{\dd x}{\dd\theta}=-\sin\theta,
\end{equation*}
se tiene
\begin{equation*}
\frac{\dd\Theta}{\dd\theta}
    =
    \frac{\dd y}{\dd x}\frac{\dd x}{\dd\theta}
    =
    -\sin\theta\frac{\dd y}{\dd x}.
\end{equation*}
Entonces
\begin{equation*}
\sin\theta\frac{\dd\Theta}{\dd\theta}
    =
    -\sin^2\theta\frac{\dd y}{\dd x}
    =
    -(1-x^2)\frac{\dd y}{\dd x}.
\end{equation*}
Derivando respecto de $\theta$,
\begin{align*}
\frac{\dd}{\dd\theta}
    \left(\sin\theta\frac{\dd\Theta}{\dd\theta}\right)
    &=
    \frac{\dd}{\dd\theta}
    \left[-(1-x^2)\frac{\dd y}{\dd x}\right]
    \notag\\
    &=
    \frac{\dd}{\dd x}
    \left[-(1-x^2)\frac{\dd y}{\dd x}\right]
    \frac{\dd x}{\dd\theta}
    \notag\\
    &=
    \sin\theta\frac{\dd}{\dd x}
    \left[(1-x^2)\frac{\dd y}{\dd x}\right].
\end{align*}
Por lo tanto,
\begin{equation*}
\frac{1}{\sin\theta}\frac{\dd}{\dd\theta}
    \left(\sin\theta\frac{\dd\Theta}{\dd\theta}\right)
    =
    \frac{\dd}{\dd x}
    \left[(1-x^2)\frac{\dd y}{\dd x}\right].
\end{equation*}
Además, como $\sin^2\theta=1-x^2$, la ecuación \eqref{eq:polar_antes_cambio} se transforma en
\begin{equation*}
\frac{\dd}{\dd x}
    \left[(1-x^2)\frac{\dd y}{\dd x}\right]
    +
    \left(Q-\frac{m^2}{1-x^2}\right)y=0,
    \qquad x\in[-1,1].
\end{equation*}
Esta es la ecuación asociada de Legendre. Al exigir soluciones regulares en los extremos $x=\pm1$, se selecciona
\begin{equation*}
Q=\ell(\ell+1),
    \qquad
    \ell=0,1,2,\ldots,
    \qquad
    |m|\leq \ell.
\end{equation*}
En ese caso, las soluciones angulares se expresan en términos de funciones asociadas de Legendre $P_\ell^m(\cos\theta)$ y la dependencia azimutal $\e^{\ii m\phi}$ se agrupa naturalmente en los armónicos esféricos $Y_\ell^m(\theta,\phi)$.

\problema{7}{Caso de simetría axial}

Si el problema físico posee simetría axial, la solución no depende de la coordenada $\phi$. Es decir,
\begin{equation*}
\Psi=\Psi(r,\theta).
\end{equation*}
Desde la formulación general, esto corresponde a tomar $\Phi(\phi)$ constante. En la ecuación azimutal
\begin{equation*}
\Phi''+m^2\Phi=0,
\end{equation*}
las soluciones constantes sólo aparecen cuando
\begin{equation*}
m=0.
\end{equation*}
Por lo tanto, la simetría axial selecciona el sector $m=0$.

Con $m=0$, la EDO polar \eqref{eq:edo_polar} se reduce a
\begin{equation*}
\frac{1}{\sin\theta}\frac{\dd}{\dd\theta}
    \left(\sin\theta\frac{\dd\Theta}{\dd\theta}\right)
    +Q\Theta=0.
\end{equation*}
Con $x=\cos\theta$, esta ecuación se transforma en
\begin{equation*}
\frac{\dd}{\dd x}\left[(1-x^2)\frac{\dd y}{\dd x}\right]+Qy=0,
\end{equation*}
o, equivalentemente,
\begin{equation*}
(1-x^2)y''-2xy'+Qy=0.
\end{equation*}
Para soluciones finitas en $x=\pm1$, se requiere
\begin{equation*}
Q=\ell(\ell+1),
    \qquad \ell=0,1,2,\ldots,
\end{equation*}
y las soluciones regulares son
\begin{equation*}
\Theta(\theta)=P_\ell(\cos\theta).
\end{equation*}

\problema{8}{Caso particular $\alpha=0$: ecuación de Laplace}

Si $\alpha=0$, la ecuación \eqref{eq:helmholtz} se reduce a la ecuación de Laplace
\begin{equation*}
\nabla^2\Psi=0.
\end{equation*}
La EDO radial \eqref{eq:edo_radial} queda
\begin{equation*}
\frac{\dd}{\dd r}\left(r^2\frac{\dd R}{\dd r}\right)-QR=0.
\end{equation*}
Expandiendo la derivada,
\begin{equation}
r^2R''+2rR'-QR=0.
    \label{eq:euler_cauchy_radial}
\end{equation}
Esta es una ecuación de Euler-Cauchy. Buscamos soluciones de la forma
\begin{equation*}
R(r)=r^s.
\end{equation*}
Entonces
\begin{equation*}
R'(r)=sr^{s-1},
    \qquad
    R''(r)=s(s-1)r^{s-2}.
\end{equation*}
Sustituyendo en \eqref{eq:euler_cauchy_radial}, resulta
\begin{align*}
r^2s(s-1)r^{s-2}+2rsr^{s-1}-Qr^s&=0,
    \notag\\
    \left[s(s-1)+2s-Q\right]r^s&=0,
    \notag\\
    \left[s(s+1)-Q\right]r^s&=0.
\end{align*}
Por lo tanto,
\begin{equation*}
s(s+1)=Q.
\end{equation*}
Para $Q=\ell(\ell+1)$, las dos soluciones son
\begin{equation*}
s=\ell,
    \qquad
    s=-(\ell+1).
\end{equation*}
Luego, la solución radial general asociada a un valor fijo de $\ell$ es
\begin{equation*}
R_\ell(r)=A_\ell r^\ell+B_\ell r^{-(\ell+1)}.
\end{equation*}
Así, en el caso de Laplace, la estructura radial queda determinada por las potencias $r^\ell$ y $r^{-(\ell+1)}$.

\problema{9}{Resultado final de la separación}

La ecuación
\begin{equation*}
\nabla^2\Psi+\alpha\Psi=0
\end{equation*}
con el ansatz
\begin{equation*}
\Psi_{\mathrm{sep}}(r,\theta,\phi)=R(r)\Theta(\theta)\Phi(\phi)
\end{equation*}
conduce al sistema
\begin{align*}
\frac{\dd}{\dd r}\left(r^2\frac{\dd R}{\dd r}\right)
    +(\alpha r^2-Q)R&=0,
    \\
    \frac{1}{\sin\theta}\frac{\dd}{\dd\theta}
    \left(\sin\theta\frac{\dd\Theta}{\dd\theta}\right)
    +
    \left(Q-\frac{m^2}{\sin^2\theta}\right)\Theta&=0,
    \\
    \frac{\dd^2\Phi}{\dd\phi^2}+m^2\Phi&=0.
\end{align*}
La periodicidad en $\phi$ impone $m\in\Z$. La regularidad angular selecciona $Q=\ell(\ell+1)$, con $\ell\in\Nzero$ y $|m|\leq\ell$.


\problema{10}{Gráficos}

\begin{figure}[H]
    \centering
    \includesvg[width=0.88\linewidth]{figuras/helmholtz_azimutal_periodicidad}
    \caption{Parte real de los modos azimutales $\Phi_m(\phi)=\e^{\ii m\phi}$ para $m=0,1,2,3$. La parte real se grafica para mostrar la periodicidad angular de forma directa. Los puntos $\phi=0$ y $\phi=2\pi$ representan la misma dirección física; por eso la función debe tomar el mismo valor en ambos extremos. Esta condición de univaluación es la que selecciona $m\in\mathbb{Z}$.}
    \label{fig:helmholtz-periodicidad}
\end{figure}

\begin{figure}[H]
    \centering
    \includesvg[width=0.88\linewidth]{figuras/helmholtz_polar_legendre_axial}
    \caption{Dependencia polar en el sector axial $m=0$. En este caso la ecuación polar se reduce a la ecuación de Legendre y las soluciones regulares son $\Theta_\ell(\theta)=P_\ell(\cos\theta)$, con $Q=\ell(\ell+1)$. El eje horizontal corresponde a $\theta/\pi$, de modo que $\theta=0$ y $\theta=\pi$ son los polos. La regularidad en esos extremos es la condición que discretiza la constante de separación angular.}
    \label{fig:helmholtz-legendre-axial}
\end{figure}

\begin{figure}[H]
    \centering
    \includesvg[width=0.88\linewidth]{figuras/helmholtz_radial_laplace_alpha_cero}
    \caption{Ramas radiales independientes para el caso $\alpha=0$, donde la ecuación de Helmholtz se reduce a la ecuación de Laplace. Para $Q=\ell(\ell+1)$ aparecen las soluciones $r^\ell$ y $r^{-(\ell+1)}$. Las curvas fueron normalizadas sólo para compararlas en una misma escala vertical; esa normalización no cambia la dependencia radial. En problemas interiores suele conservarse la rama regular $r^\ell$, mientras que en problemas exteriores suele aparecer la rama decreciente $r^{-(\ell+1)}$.}
    \label{fig:helmholtz-radial-laplace}
\end{figure}

\end{document}
